\chapter{讨论与结论}\label{sec:discussion}

两条财富路径到达同一终点后，参考点仍可能不同；投资者下一期面对的也就不是同一个CPT目标。在财富、持仓和可选投资机会相同的条件下，本文的路径比较确实得到不同的风险资产配置。由此，参考点既记录已发生的投资经历，也参与决定下一期策略所优化的目标。针对每一期给定状态后的目标，本文构造严格无偏且方差有限的梯度估计量；当市场与账户状态推动目标继续移动时，策略再依据估计梯度更新，并用投影残差和动态局部遗憾衡量它能否跟上当期的一阶条件。

固定参考点的CPT强化学习已经研究了非线性概率加权下的策略梯度优化\citep{prashanth2016cumulative,lepel2026policygradients}；行为金融研究则表明，已实现的收益与损失会改变投资者评价后续结果的基准\citep{arkes2008reference,he2019realization,qin2022realization,gao2024multiperiod}。本文把随账户路径更新的参考点纳入连续策略学习，使这两条研究线索作用于同一个投资过程。有限轨迹预算下的梯度估计为逐期更新提供了一阶信息；在缓变市场中，持续更新比冻结参数积累更少的跟踪误差。参考点适应消融还显示，即使参考点保持静止，在线更新仍然有效，而收益主导的非对称适应可以进一步增强这种优势。参考点影响目标的移动方式，参数更新则决定策略如何响应这种移动。

按时间顺序进行的A股回放将比较从一阶跟踪延伸到真实价格路径上的投资后果。DRCPT-PG与对称参考点方法的平均收益接近，但平均Sharpe比率较高，最大回撤、换手和交易成本较低；收益相近时，参考点机制的差异更多表现在持仓风险和调仓强度上。真实投资者的持股决策又提供了行为方面的证据：在相同样本、控制变量和测试时段下，非对称动态参考点提高了样本外减持预测能力，按这一参考点划分的盈亏状态也更贴近实际减持差异。这些差别来自对已观察行为的解释与预测，不涉及参考点对卖出行为的因果识别。以历史行为校准的模拟投资者中，各方法的采纳率相近，采纳后相对“不调仓”的模型前景价值却有所不同。投资者在期初愿意接受一份建议，与五日之后如何评价它，是两个相继发生但不等同的问题。

现有跟踪结果建立在目标缓慢变化的市场路径上。面对突发且持续的冲击，目标变化与可用轨迹预算都可能更快增加；如何在这种情况下控制一阶跟踪误差，是下一步需要解决的问题。模拟交互中的投资者类型来自历史行为校准，条件满意度则是模型计算的五日前景价值差值。若要进一步研究投资者的真实感受，需要在持续交互中同时观察参考点变化、推荐采纳、实际投资结果和主观评价。本文说明了历史经历如何进入可估计、可跟踪的投资组合目标；把这一目标与可直接观察的人类体验连接起来，仍需相应的行为数据。
